\newpage
\section{Appendix E}
This section presents the full, unannotated code used for the extraction of data. Note that for 1911 and 1914, the miles code used 'Track and Equipment' instead of 'Plant and Equipment'. However, as the code is elsewise identical, it has not been repeated in its own section. '---' is used to represent em-dashes

The section also gives an annotated copy of the code used to fit S-curve models.

\subsection{Extraction of entries from 1894-1914}
\begin{lstlisting}[language=Python]
import pandas as pd
import re
import matplotlib.pyplot as plt
import numpy as np
from re import search
import string
from fuzzywuzzy import fuzz
from fuzzywuzzy import process
import csv

def read_file(year):
    out = None
    with open(
        f'{year} TXT.txt', #text file name
        encoding="utf8"
    ) as f:
        out = pd.DataFrame(f)
        out['Info'] = out[0].apply(lambda x: x.strip())
        out['Year'] = year
        out = out.drop(columns=[0])
    return out

def hasNumbers(inputString):
    return bool(re.search(r'\d', inputString))   

def capratio(bits):
    start = 0
    ratio = 0
    translator = bits.maketrans('', '', string.punctuation)
    ring = bits.translate(translator)
    ringing = ring.replace(' ', '')
    
    for i in range(len(bits)):
        if bits[i].isupper() is True:
            start = start + 1
    if len(ringing) > 0:
        ratio = start/len(ringing)
    
    return ratio
    
def getcities(Info):
if capratio(Info) > 0.7 and len(Info) < 30:
    return Info
elif len(Info) in range(7,25) and process.extractOne(Info.lower(), cities.str.lower(), scorer=fuzz.ratio)[1] > 80 and hasNumbers(Info) is False:
    if abs(len(process.extractOne(Info.lower(), cities.str.lower(), scorer=fuzz.ratio)[0]) - len(Info.lower())) < 4:
        return Info
else:
    return 999
    
def city2(CityInfo): 
    if CityInfo != 999 and CityInfo is not None:
        section = CityInfo
        if process.extractOne(section.lower(), cities, scorer=fuzz.ratio)[1] > 80:
            return process.extractOne(section.lower(), cities, scorer=fuzz.ratio)[0]
        elif process.extractOne(section.lower(), cities, scorer=fuzz.partial_ratio)[1] > 98 and process.extractOne(section.lower(), cities, scorer=fuzz.ratio)[1] > 50:
            return process.extractOne(section.lower(), cities, scorer=fuzz.partial_ratio)[0]
        else:
            return np.nan
    else:
        return np.nan

def getcompanies(Info, CleanCity, CompIndex, Truth):
    try:
        if len(Info) > 5:
            if len(re.findall('Funded|leased|Leased|See|Deficit|Interest', Info[0:40])) == 0 and '....' not in Info:
                if len(re.findall('Charter|Incorp|Succe|consolidation|Connects|This|connect|Organize|organize|Owns|succe|Annual|Formed|nchise', Info)) > 0 or Truth:
                    for x in CompIndex:
                        if fuzz.ratio(x[1], CleanCity) > 90:
                            if x[0] in Info[0:70]:
                                return x[0]
                            elif fuzz.partial_ratio(x[0].lower(), Info[0:70].lower()) > 90 and len(Info) > 11:
                                return x[0]
                            elif 'Population' in Info:
                                if fuzz.partial_ratio(x[0].lower(), Info[0:140].lower()) > 90:
                                    return x[0]
                            else:
                                pass
                        else:
                            pass
                else:
                    pass
            else:
                pass
        else:
            pass
    except:
        pass
    
def getcompanies2(Info, CleanCity, CurrentList, CompIndex, Truth):
    try:
        if len(Info) > 5:
            if CurrentList is None:
                if len(re.findall('Funded|leased|Leased|See|Deficit|Interest', Info[0:40])) == 0 and '....' not in Info:
                    if len(re.findall('Charter|Incorp|Succe|consolidation|Connects|This|connect|Organize|organize|Owns|succe|Annual|Formed|nchise', Info)) > 0:
                        for x in CompIndex:
                            if fuzz.ratio(x[1], CleanCity) > 80:
                                if fuzz.partial_ratio(x[0].lower(), Info[0:70].lower()) > 80 and len(Info) > 11: 
                                    return x[0]
                                elif 'Population' in Info:
                                    if fuzz.partial_ratio(x[0].lower(), Info[0:140].lower()) > 80:
                                        return x[0]
                                else:
                                    pass
                            else:
                                pass
                        
                    else:
                        pass
                else:
                    pass
            else:
                return CurrentList
        else:
            pass
    except:
        pass
        
def getmiles(Info): 
    if 'Contemplated' not in Info and 'Proposed' not in Info and 'under construction' not in Info[0:40] and 'power plant' not in Info[0:40]:
        if fuzz.partial_ratio(Info, 'Plant and Equipment') > 80 and len(Info) > 20 and 'total' in Info:
            try:
                split_string = Info.split('total')[1]
                mile = (re.findall("[+-]?\d+\.+\d+|\d+", split_string))[0]
                return mile
            except:
                return 'No'
        elif fuzz.partial_ratio(Info, 'Plant and Equipment') > 80 and len(Info) > 20:
            try:
                split_string = Info.split('ment')[1]        
                mile = (re.findall("[+-]?\d+\.+\d+|\d+", split_string))[0]
                return mile
            except:
                return 'No'
        elif 'Miles of track' in Info[0:60] and 'total' in Info:
            split_string = Info.split('total')[1]
            mile = (re.findall("[+-]?\d+\.+\d+|\d+", split_string))[0]
            return mile
        elif 'Miles of track' in Info[0:60]:
            split_string = Info.split('of track')[1]
            try:
                mile = (re.findall("[+-]?\d+\.+\d+|\d+", split_string))[0]
                return mile
            except:
                return 'No'
        elif 'Miles of Track' in Info[0:60]:
            split_string = Info.split('of track')[1]
            try:
                mile = (re.findall("[+-]?\d+\.+\d+|\d+", split_string))[0]
                return mile
            except:
                return 'No'
        else:
            return 'No'
    else:
        return 'No'
        
def getstate(Companies, RealCities, CompIndex):
    for x in CompIndex:
        if Companies == x[0] and RealCities == x[1]:
            return x[2]
            
years = [1894] + list(range(1898, 1915))
for year in years:
    try:
        comps = pd.read_csv(f'{year}.csv')
        comps = comps.drop(['Unnamed: 0'], axis = 1)

        companies = comps['Company']
        places = comps['Place']
        cities = comps['Place'].str.split(',').str[0]
        states = comps['Place'].str.split(',').str[-1]

        index = list(zip(companies, cities, states))

        df = read_file(year)

        value = len(df)
        start = 0
        end = 0

        for i in range(value):
            if 'START OF RECORDS' in df.iloc[i,0]:
                start = i
            elif 'END OF RECORDS' in df.iloc[i,0]:
                end = i+1

        df = df.drop(df.index[end:value])
        df = df.drop(df.index[0:start])
        df = df.reset_index(drop=True)

        df['Info'] = df['Info'].str.replace('\t', ' ')
        df['Info'] = df['Info'].str.replace('. 0.', '. O.')
        df['Info'] = df['Info'].str.replace(', 0.', ', O.')
        df = df[df['Info'].map(len) > 4]

        df['Info'].replace('', np.nan, inplace=True)
        df.dropna(subset=['Info'], inplace=True)
        df = df.reset_index(drop=True)

        unwanted = []

        for x in df['Info']:
            if 'MANUAL' in x or 'McGRAW' in x or 'AMERICAN STREET' in x or 'INVESTMENTS' in x or 'UNIVERSITY' in x:
                unwanted.append(x)

        unwanted = list(set(unwanted))

        df['Unwanted'] = df['Info'].isin(unwanted)

        listA= df[df['Unwanted'] == True].index.tolist()

        listB= []

        for x in listA:
            try:
                if capratio(df['Info'][x+1]) < 0.7 and len(df['Info'][x+1]) < 20:
                   listB.append(x+1)
                if capratio(df['Info'][x-1]) < 0.7 and len(df['Info'][x-1]) < 20:
                   listB.append(x-1)
                if capratio(df['Info'][x+2]) < 0.7 and len(df['Info'][x+2]) < 20:
                   listB.append(x+2)
                if capratio(df['Info'][x-2]) < 0.7 and len(df['Info'][x-2]) < 20:
                   listB.append(x-2)
                if capratio(df['Info'][x+3]) < 0.7 and len(df['Info'][x+3]) < 20:
                   listB.append(x+3)
                if capratio(df['Info'][x-3]) < 0.7 and len(df['Info'][x-3]) < 20:
                   listB.append(x-3)
                if capratio(df['Info'][x+4]) < 0.7 and len(df['Info'][x+4]) < 20:
                   listB.append(x+4)
                if capratio(df['Info'][x-4]) < 0.7 and len(df['Info'][x-4]) < 20:
                   listB.append(x-4)
            except:
                pass

        listC = listA+listB
        indexes_to_keep = list(set(list(range(len(df)))) - set(listC))
        df = df.take(indexes_to_keep)   
        df = df.reset_index(drop=True)
        
        df['Truth'] = df['Info'].str.contains('Capital Stock|Funded Debt')
        for x in range(len(df['Truth'])):
            if df['Truth'][x]:
                df['Truth'][x-1] = True
                df['Truth'][x-2] = True
    
        df['Cities'] = df['Info'].apply(getcities)
        df['RealCities'] = df['Cities'].apply(city2)
        df['RealCities'] = df['RealCities'].ffill()

        df['Miles'] = df['Info'].apply(getmiles)

        df['Company'] = df[1:len(df)].apply(lambda x: getcompanies(x.Info, x.RealCities, index, x.Truth), axis=1)
        df['Company'] = df[1:len(df)].apply(lambda x: getcompanies2(x.Info, x.RealCities, x.Company, index, x.Truth), axis=1)

        for i in range(len(df)):
            try:
                if df['Company'][i] is not None and df['Company'][i+1] is None and df['RealCities'][i] == df['RealCities'][i+1]:
                    df['Company'][i+1] = df['Company'][i]
            except:
                pass

        df['State'] = df.apply(lambda x: getstate(x.Company, x.RealCities, index), axis=1)

        extract = []

        for x in range(len(df)):
            if df['Miles'][x] != 'No':
                extract.append(x)
        results = df.take(extract)
        results = results.drop(columns=['Info', 'Cities', 'Unwanted'])
        display(results)
    except:
        pass
            
\end{lstlisting}

\subsection{Extraction of entries from 1917 onwards}
\begin{lstlisting}
import pandas as pd
import re
import matplotlib.pyplot as plt
import numpy as np
import us
from re import search
import string
from fuzzywuzzy import fuzz
from fuzzywuzzy import process
import csv

def read_file(year):
    out = None
    with open(
        f'{year} TXT.txt', #text file name
        encoding="utf8"
    ) as f:
        out = pd.DataFrame(f)
        out['Info'] = out[0].apply(lambda x: x.strip())
        out['Year'] = year
        out = out.drop(columns=[0])
    return out

def hasNumbers(inputString):
    return bool(re.search(r'\d', inputString))   

def capratio(bits):
    bits = ''.join((x for x in bits if not x.isdigit()))
    start = 0
    ratio = 0
    translator = bits.maketrans('', '', string.punctuation)
    ring = bits.translate(translator)
    ringing = ring.replace(' ', '')
    
    for i in range(len(bits)):
        if bits[i].isupper() is True:
            start = start + 1
    if len(ringing) > 0:
        ratio = start/len(ringing)
    
    return ratio
    
def get_state(x):
    if x in states and hasNumbers(x) is False:
        return x
    else:
        return np.nan
    
def get_city(x):
    if capratio(x) > 0.7 and len(x) > 5 and hasNumbers(x) is True and 'STA.' not in x and 'EQUIP' not in x:
        x.replace('8', 'S')
        x.replace('5', 'S')
        if 'ST.' not in x and 'FT.' and 'MT.' not in x:
            try:
                return x.split('.')[0]
            except:
                return np.nan
        else:
            splits = x.split('.')
            try:
                if hasNumbers(splits[1]) is False:
                    return splits[0] + splits[1]
                else:
                    return splits[0]
            except:
                return np.nan
    else:
        return np.nan

def get_company(x):
    if '---' in x and hasNumbers(x[0:2]) is True:
        if 'Co.' in x:  #one for split at emdash and Co.
            try:
                compstring = x.split('---')[1]
                if len(compstring) > 2:
                    company = compstring.split('Co.')[0] + 'Co.'
                    return company
                else:
                    compstring = x.split('---')[0]
                    company = ''.join((x for x in compstring if not x.isdigit()))
                    return company
            except:
                return np.nan
        elif len(x.split('---')) > 2:
            return x.split('---')[1]
        else:
            try:
                return x.split('---')[1]
            except:
                return np.nan
    else:
        return np.nan
        
def get_mile(x):
    x = x.lower()
    if len(x) > 20:
        if 'miles owned' in x: #simplest
            splitter = 'miles owned'
            try:
                str_to_check = x.split(splitter)[0]
                num_list = re.findall("[+-]?\d+\.+\d+|[+-]?\d+\,+\d+|\d+ \d+|\d+", str_to_check)
                mile = num_list[len(num_list)-1]
                return mile
            except:
                return 'No'
        elif '4-8' in x:
            splits = x.split('4-8')
            checkpoint = ''.join('' if x is splits[len(splits)-1] else x for x in splits)
            longs = len(checkpoint) - 1
            if longs > 180:
                shorts = longs - 180
            else:
                shorts = 0
            alist = checkpoint[shorts:longs].split(' ')
            alist = [x for x in alist if len(x) > 4]
            if fuzz.partial_ratio('miles', x[shorts:longs]) > 75: #miles is in the string
                try:
                    splitter = process.extractOne('miles', alist, scorer=fuzz.partial_ratio)[0]
                    try:
                        str_to_check = x[shorts:longs].split(splitter)[0]
                        num_list = re.findall("[+-]?\d+\.+\d+|[+-]?\d+\,+\d+|\d+ \d+|\d+", str_to_check)
                        mile = num_list[len(num_list)-1]
                        return mile
                    except:
                        return 'No'
                except:
                    pass
            else:
                return 'No'
        elif '5-2' in x and '5-2000' not in x:
            splits = x.split('5-2')
            checkpoint = ''.join('' if x is splits[len(splits)-1] else x for x in splits)
            longs = len(checkpoint) - 1
            if longs > 180:
                shorts = longs - 180
            else:
                shorts = 0
            alist = checkpoint[shorts:longs].split(' ')
            alist = [x for x in alist if len(x) > 4 and x is not None]
            if fuzz.partial_ratio('miles', x[shorts:longs]) > 75: #miles is in the string
                splitter = process.extractOne('miles', alist, scorer=fuzz.partial_ratio)[0]
                try:
                    str_to_check = x[shorts:longs].split(splitter)[0]
                    num_list = re.findall("[+-]?\d+\.+\d+|[+-]?\d+\,+\d+|\d+ \d+|\d+", str_to_check)
                    mile = num_list[len(num_list)-1]
                    return mile
                except:
                    return 'No'
            else:
                return 'No'
        else:
            longs = len(x) - 1
            if longs > 180:
                shorts = longs - 180
            else:
                shorts = 0
            if 'miles' in x[shorts:longs]:
                str_to_check = x[shorts:longs].split('miles')[0]
                try:
                    num_list = re.findall("[+-]?\d+\.+\d+|[+-]?\d+\,+\d+|\d+ \d+|\d+", str_to_check)
                    mile = num_list[len(num_list)-1]
                    return mile
                except:
                    return 'No'
            else:
                return 'No'            
    else:
        return 'No'
        
states1 = [str(i).upper() for i in us.states.STATES]
states2 = [str(i).upper() + '.' for i in us.states.STATES]
states = states1 + states2

years = list(range(1917, 1927))

for year in years:
    df = read_file(year)
    
    value = len(df)
    start = 0
    end = 0
    
    for i in range(value):
            if 'START OF RECORDS' in df['Info'][i]:
                start = i
            elif 'END OF RECORDS' in df['Info'][i]:
                end = i+1

    df = df.drop(df.index[end:value])
    df = df.drop(df.index[0:start])
    df = df.reset_index(drop=True)
    
    df['Info'] = df['Info'].str.replace('. 0.', '. O.')
    df['Info'] = df['Info'].str.replace(', 0.', ', O.')
    df['Info'] = df['Info'].str.replace('\t', ' ')
    df = df[df["Info"].str.contains("UNIVERSITY")==False]
    df = df.reset_index(drop=True)
    
    df['State'] = df['Info'].apply(get_state)
    df['State'] = df['State'].ffill()
    
    df['City'] = df['Info'].apply(get_city)
    df['City'] = df['City'].str.replace('8', 'S')
    df['City'] = df['City'].ffill()

    df['Miles'] = df['Info'].apply(get_mile)
    df['Company'] = df['Info'].apply(get_company)
    df['Company'] = df['Company'].ffill()
    
    extract = []
    
    for x in range(len(df)):
            if df['Miles'][x] != 'No':
                extract.append(x)
    results = df.take(extract)
    results = results.drop(columns=['Info'])
    display(results)
\end{lstlisting}

\subsection{Fitting the S-curve model}
\begin{lstlisting}[language=Python]
import pandas as pd
import matplotlib.pyplot as plt
import numpy as np
import csv
import math
import statsmodels.api as sm
import statsmodels.formula.api as smf
import us
from fuzzywuzzy import fuzz

states = [str(i) for i in us.states.STATES] + ['D. C']
df = pd.read_csv('data.csv')
df['Year'] = df['Year'].astype(int)

xval = []
yval = []

years = [1894] + list(range(1898,1912)) + [1914] + list(range(1917, 1927))

place = [] #to record state
k_optimal = [] #to record best k value
params_intercept = [] #to record intercept
params_slope = [] #to record slope
params_rsquared_adj = [] #to record adj rsquared
params_pvalue = [] #to record pvalue
params_stderr_int = [] #to record intercept std error
params_stderr_slp = [] #to record slope std error
params_stderr_reg = [] #to record regression std error
time_inflex = [] #to record t0 
observes = [] #to record sample size
incdec = [] #to record if increase or decrease

for state in states:
    
        for year in years:

            #gathering data
            year_sum = 0

            for i in range(len(df)):
                if df['State'][i] == state and df['Year'][i] == year:
                    year_sum = year_sum + df['Miles'][i]

            yval.append(year_sum)
            xval.append(year)

        #finding s-curve
        k_loc = yval.index(max(yval)) #find the year with max track length

        if k_loc > len(yval) - 3: #it's very close to 1926, no need to make 2 separate s-curves
            k_start = int(max(yval)) + 1
            k_range = range(k_start + 1, round(k_start*1.5), 1) #generates a range of k-values to test, up to 150% of maximum

            reg_calc = []
            rsq_value = -1
            best_k = 0
            best_reg_calc = []

            for value in k_range:
                for i in range(len(yval)):
                    if yval[i] != 0: #if not a 0 mile year, use the equation ln(miles/k minus miles)
                        from_ln = np.log(yval[i]/(value - yval[i]))
                    else:
                        from_ln = -2 #low discouraging value for regression

                    reg_calc.append(from_ln) #save the value

                #once all reg_calc for that value found, find correlation coefficient, rsq
                correlation_matrix = np.corrcoef(xval, reg_calc)
                correlation_xy = correlation_matrix[0,1]
                r_squared = correlation_xy**2

                if r_squared > rsq_value: #if it's a better rsq than before, store it:
                    best_k = value
                    best_reg_calc = reg_calc

                    slope_value, slt0 = np.polyfit(xval, reg_calc, deg=1)
                    time_0 = -1 * slt0/slope_value

                reg_calc = [] #clear the reg_calc list so the loop can rerun for next k-value


            #once optimum k-value acquired:

            s_t = []
            for i in range(len(best_reg_calc)):
                power = xval[i] - time_0
                denom = 1 + math.exp(-1 * slope_value * power)
                total = best_k/denom
                s_t.append(total) #the k/1+e^-b(t-t0) equation

            plt.plot(xval, s_t, "tab:orange", label="S-curve")
            plt.plot(xval, yval, "tab:blue", label="McGraw Hill Data")
            plt.legend()
            plt.title(state)
            plt.xlim(1890, 1930)
            plt.xlabel('Year')
            plt.ylabel('Miles of track')
            plt.show()
            plt.clf()

            data2 = pd.DataFrame({'x':xval, 'y':best_reg_calc})
            model = smf.ols('y~x', data=data2).fit()

            #recording of stats
            place.append(state)
            observes.append(len(xval))
            time_inflex.append(time_0)
            k_optimal.append(best_k)
            params_intercept.append(model.params[0])
            params_slope.append(model.params[1])
            params_rsquared_adj.append(model.rsquared_adj) 
            params_pvalue.append(model.f_pvalue)
            params_stderr_int.append(model.bse[0])
            params_stderr_slp.append(model.bse[1])
            params_stderr_reg.append(model.scale**0.5)
            incdec.append('Inc')

            xval = []
            yval = []

        else: #if there does need to be two s-curves because the peak is in the middle
            inc_years = xval[:k_loc+1] #split into years with increasing values and years with decreasing values
            dec_years = xval[k_loc:] 

            inc_values = yval[:k_loc+1] #split into increasing values and decreasing values
            dec_values = yval[k_loc:]
            
            ##for regression with non-zero limits only: remove these two lines if regression with zero limit
            dec_years.append(1940)
            dec_values.append(max(yval)/3)

            #for INCREASING years
            k_start = int(max(yval)) + 1
            k_range = range(k_start + 1, k_start + 10, 1) #the range of k-values to test - max to max+10 (since known max)

            reg_calc = []
            rsq_value = -1
            best_k = 0
            best_reg_calc = []

            for value in k_range:
                for i in range(len(inc_values)):
                    if inc_values[i] != 0:
                        from_ln = np.log(inc_values[i]/(value-inc_values[i])) #the ln(miles/k minus miles) equation
                    else:
                        from_ln = -2
                    
                    reg_calc.append(from_ln)

                correlation_matrix = np.corrcoef(inc_years, reg_calc) #find the correlation coefficient and r squared value
                correlation_xy = correlation_matrix[0,1]
                r_squared = correlation_xy**2

                if r_squared > rsq_value: #if it's a better rsq than before, store it:
                    rsq_value = r_squared
                    best_k = value
                    best_reg_calc = reg_calc

                    slope_value, slt0 = np.polyfit(inc_years, reg_calc, deg=1)
                    time_0 = -1 * slt0/slope_value

                reg_calc = [] #clear the equation to run loop again

            inc_s_t = [] #once the best reg_calc has been found:

            for i in range(len(best_reg_calc)):
                power = inc_years[i] - time_0
                denom = 1 + math.exp(-1 * slope_value * power)
                total = best_k/denom
                inc_s_t.append(total)

            data2 = pd.DataFrame({'x':inc_years, 'y':best_reg_calc})
            model = smf.ols('y~x', data=data2).fit()

            place.append(state)
            observes.append(len(inc_years))
            time_inflex.append(time_0)
            k_optimal.append(best_k)
            params_intercept.append(model.params[0])
            params_slope.append(model.params[1])
            params_rsquared_adj.append(model.rsquared_adj) 
            params_pvalue.append(model.f_pvalue)
            params_stderr_int.append(model.bse[0])
            params_stderr_slp.append(model.bse[1])
            params_stderr_reg.append(model.scale**0.5)
            incdec.append('Inc')

            #for DECREASING years

            dec_reg_calc = []

            for i in range(len(dec_values)):
                from_ln = np.log(dec_values[i]/(best_k-dec_values[i])) #the ln(miles/k minus miles) equation
                dec_reg_calc.append(from_ln)

            correlation_matrix = np.corrcoef(dec_years, dec_reg_calc) #find the correlation coefficient and r squared value
            correlation_xy = correlation_matrix[0,1]
            r_squared = correlation_xy**2

            slope_value, slt0 = np.polyfit(dec_years, dec_reg_calc, deg=1)
            time_0 = -1 * slt0/slope_value

            dec_s_t = [] 

            for i in range(len(dec_reg_calc)):
                power = dec_years[i] - time_0
                denom = 1 + math.exp(-1 * slope_value * power)
                total = best_k/denom
                dec_s_t.append(total)

            data3 = pd.DataFrame({'x':dec_years, 'y':dec_reg_calc})
            model = smf.ols('y~x', data=data3).fit()

            place.append(state)
            observes.append(len(dec_years))
            time_inflex.append(time_0)
            k_optimal.append(best_k)
            params_intercept.append(model.params[0])
            params_slope.append(model.params[1])
            params_rsquared_adj.append(model.rsquared_adj) 
            params_pvalue.append(model.f_pvalue)
            params_stderr_int.append(model.bse[0])
            params_stderr_slp.append(model.bse[1])
            params_stderr_reg.append(model.scale**0.5)
            incdec.append('Dec')
            
            ##for non-zero regression ONLY
            dec_years = dec_years[:-1]
            dec_values = dec_values[:-1]

            plt.plot(xval, yval, "tab:blue", label="McGraw Hill Data")
            plt.plot(inc_years, inc_s_t, "tab:orange", label="Increasing S-Curve")
            plt.plot(dec_years, dec_s_t, "tab:pink", label="Decreasing S-Curve")
            plt.legend()
            plt.title(state)
            plt.xlabel('Year')
            plt.ylabel('Miles of track')
            plt.show()

            xval = []
            yval = []

result = pd.DataFrame({
        'State' : place,
        'Sample Size': observes,
        'K-value' : k_optimal,
        'Time Zero' : time_inflex,
        'Intercept' : params_intercept,
        'Slope' : params_slope,
        'RSQ-adj' : params_rsquared_adj,
        'Std error' : params_stderr_reg,
        'P-value' : params_pvalue,
        'Intercept error' : params_stderr_int,
        'Slope error' : params_stderr_slp,
        'IncDec' : incdec
    })

display(result)
\end{lstlisting}